%%%PAGE 9 rot=0 legibility=good summary=功的计算（恒力/变力三法）、功率；人对自己做功的能量转化表；非质点（绳/杆）做功；弹簧与能量守恒；机车启动（恒P、恒a）公式与图像
%%%BLOCK id=009-01 ch=06 sec=功·功的计算 kind=knowledge alt=none cont=none y=0.03-0.30
$W=Fs\cos\theta$\quad $s$为对地位移 % 公式下方有红笔下划线

功\quad $-8\unit{J}>5\unit{J}$

能\quad $-8\unit{J}<5\unit{J}$

定力做功\quad $W=\left\{\begin{array}{l}\bar{F}s=\dfrac{F_0+F_1}{2}\,s\\[4pt] \text{$F$-$s$图面积}\\[4pt] \text{动能定理}\end{array}\right.$ % CHECK: 原稿首字形似“定”（“定力做功”），括号内三种方法属变力求功，疑为“变力做功”，照原样转录

变力做功\quad 动能定理
%%%END
%%%BLOCK id=009-02 ch=06 sec=功率 kind=formula alt=none cont=none y=0.30-0.34
平均功率\quad $\bar{P}=\dfrac{W}{t}$

瞬时功率\quad $P=Fv\cos\alpha$
%%%END
%%%BLOCK id=009-03 ch=06 sec=功·人对自己做功 kind=knowledge alt=none cont=none y=0.33-0.42
人对自己做功（外界不对人做功）

不知去向的能量（转化为内能）

\begin{tabular}{@{}lll@{}}
加速： & $f_{\text{静}}$与$v$同向 & 生物质能$\to$机械能 \\
匀速 & $f_{\text{静}}=0$ & 生物质能\textsubscript{\unclear{（原地起跳）}}$\to$内能 \\
减速 & $f_{\text{静}}$与$v$反向 & 机械能$+$生物质能$\to$内\unclear{能} \\
\end{tabular}
% 原稿此表位于页面右侧；“加速/匀速/减速”三词带下划线；“匀速”行“生物质能”下方有一行小字（形似“（原地起跳）”）；末行箭头后的字被页边截断
%%%END
%%%BLOCK id=009-04 ch=06 sec=功·非质点对象做功 kind=method alt=none cont=none y=0.43-0.56
非质点对象做功，有质量的绳/杆

\hspace{4em}分段$+$找重心，$W=m_1g\Delta h_1+m_2g\Delta h_2$

%FIG p009_a 0.065 0.498 0.19 0.553
\notefig[0.25]{figures/p009_a.png}
%%%END
%%%BLOCK id=009-05 ch=06 sec=弹簧模型·能量守恒 kind=method alt=07 cont=none y=0.56-0.67
过程写动能\qquad 过程写动量

弹簧\quad $\left\{\begin{array}{l}\pm W_{\text{弹}}=E_{p\text{初}}-E_{p\text{末}}\Rightarrow\text{过程写动能/能量}\\[4pt] F=k\Delta x\to\text{状态写牛二}\end{array}\right.$

能量守恒\quad $\Delta E_{\text{增}}=\Delta E_{\text{减}}$（定量）\quad $E_1+\cdots E_n=$定值（定性）
%%%END
%%%BLOCK id=009-06 ch=06 sec=机车启动 kind=method alt=03 cont=none y=0.69-1.00
状态写牛二\quad $F_{\text{牵}}-f=ma$

牵引力/功率\quad $P=F_{\text{牵}}v=fv_{\text{匀}}$

匀速速率\quad $V_{\text{均}}=\dfrac{P_{\text{额}}}{f}$ % CHECK: 速度下标手写形似“均”，按物理应为“匀”，照原样转录

%FIG p009_b 0.012 0.803 0.865 0.999
\notefig[0.85]{figures/p009_b.png}
% 图 p009_b 包含：左半“恒P”（P-t、v-t、F牵-1/v、a-1/v四图）与右半“恒a”（P-t、v-t、F牵-1/v、a-1/v四图）及其中的标注；图下方两条独立公式如下

$\dfrac{P_0}{v}-f=ma$\qquad $v_1=\dfrac{P_0}{ma+f}$
%%%END
%%%PAGEEND 9
